\section{Lecture 22 (March 7, 2025, est.): Ostrik's theorem}
\label{lec22}
\noindent\textit{Scribe: Guillermo Sanmarco}

Let $k$ be algebraically closed of characteristic $p>0$.
We finish the proof that every symmetric fusion category over $k$ admits
a symmetric tensor functor to $\mathsf{Ver}_p$. The argument follows the
Frobenius functor and its fusion image, using the lemmas of the preceding
lecture.

\subsection{Separating the two factors}
Assume first that $p>3$, that $\mathcal C$ is Frobenius injective, and that
its Frobenius type is $\mathsf{Ver}_p^+$ or $\mathsf{Ver}_p$.
Define $\mathcal C_s$ to be the full fusion subcategory whose simple
objects satisfy
\[
 \mathrm{Fr}(X)=Y\boxtimes\mathbf1
 \quad\text{or}\quad
 \mathrm{Fr}(X)=Y\boxtimes\varepsilon,
 \qquad \varepsilon=L_{p-1}.
\]
This is indeed a fusion subcategory: it is the inverse image of
$\mathcal C^{(1)}\boxtimes\SVec$ under $\mathrm{Fr}$.
By naturality, its own Frobenius functor is the restriction of the one on
$\mathcal C$. Thus it is Frobenius injective and has type $\Vect$ or
$\SVec$. Lemmas~\ref{lec21-type-vec} and \ref{lec21-type-svec} imply
\begin{equation}\label{lec22-cs-nondegenerate}
 \dim(\mathcal C_s)\ne0.
\end{equation}

\begin{lemma}\label{lec22-image-factorization}
Under these assumptions,
\[
 \mathrm{Fr}(\mathcal C)
 =\langle\mathrm{Fr}(\mathcal C_s),
            \mathbf1\boxtimes\mathsf{Ver}_p^+\rangle
 \simeq\mathcal C_s\boxtimes\mathsf{Ver}_p^+
\]
as symmetric fusion categories.
\end{lemma}
\begin{proof}
Both subcategories on the right belong to the Frobenius image, the second
by Lemma~\ref{lec21-ver-plus-in-image}. For the reverse inclusion take a
simple $X\in\mathcal C$. Its image is simple, so write
$\mathrm{Fr}(X)=T\boxtimes L_r$. Put
\[
 \delta_r=
 \begin{cases}\mathbf1,&r\text{ odd},\\
 \varepsilon,&r\text{ even}.
 \end{cases}
\]
Then $L_r\otimes\delta_r$ belongs to $\mathsf{Ver}_p^+$.
Full faithfulness and semisimplicity imply that every object in the
fusion image is the image of an object: a direct summand is specified by
an idempotent, which lifts through a fully faithful functor and splits in
the source. Therefore there is a simple $Y\in\mathcal C$ with
\[
 \mathrm{Fr}(Y)=\mathbf1\boxtimes(L_r\otimes\delta_r).
\]
Now $\mathrm{Fr}(X\otimes Y)$ has the simple summand
$T\boxtimes\delta_r$, because $L_r$ is self-dual and
$\mathbf1$ is a summand of $L_r\otimes L_r$. Lift this summand to a
simple $Z$ of $X\otimes Y$. It satisfies
$\mathrm{Fr}(Z)=T\boxtimes\delta_r$, hence $Z\in\mathcal C_s$, and
\[
 \mathrm{Fr}(X)
 \simeq\mathrm{Fr}(Z)\otimes
       \bigl(\mathbf1\boxtimes(L_r\otimes\delta_r)\bigr).
\]
This proves generation.

It remains to justify that generation is a Deligne product. The two
subcategories have only the unit as a common simple: a simple in
$\mathrm{Fr}(\mathcal C_s)$ has second factor $\mathbf1$ or
$\varepsilon$, whereas a simple in
$\mathbf1\boxtimes\mathsf{Ver}_p^+$ has first factor $\mathbf1$ and
odd-indexed second factor. The only common possibility is the unit.
In a symmetric fusion category, two full fusion subcategories
$\mathcal A,\mathcal B$ with this property have fully faithful product
functor $\mathcal A\boxtimes\mathcal B\to\langle\mathcal A,\mathcal B\rangle$.
Indeed, for simple $A_i\in\mathcal A$ and $B_i\in\mathcal B$, rigidity
identifies
\[
 \Hom(A_1\otimes B_1,A_2\otimes B_2)
 \simeq\Hom(A_2^*\otimes A_1,B_2\otimes B_1^*).
\]
Only the unit can contribute to this Hom-space; it contributes once
exactly when $A_1\simeq A_2$ and $B_1\simeq B_2$. The product functor
is therefore fully faithful and, by generation, essentially surjective.
Apply this with
$\mathcal A=\mathrm{Fr}(\mathcal C_s)$ and
$\mathcal B=\mathbf1\boxtimes\mathsf{Ver}_p^+$, and use
$\mathrm{Fr}(\mathcal C_s)\simeq\mathcal C_s$.
\end{proof}

\subsection{The theorem}
\begin{theorem}[Ostrik]\label{lec22-ostrik}
Every symmetric fusion category $\mathcal C$ over $k$ admits a $k$-linear
exact symmetric monoidal functor
\[
 \mathcal C\longrightarrow\mathsf{Ver}_p.
\]
\end{theorem}
\begin{proof}
Suppose there is a counterexample and choose one of minimal
Frobenius--Perron dimension. Such a minimum exists. To check this point,
fix a bound $M$. A fusion category of dimension at most $M$ has at most
$M$ simple objects, because every simple has Frobenius--Perron dimension
at least one. Its simple dimensions are at most $\sqrt M$, and its fusion
multiplicities satisfy
$N_{ij}^r\leq d_i d_j/d_r\leq M$. There are therefore only finitely many
possible based fusion rings below the bound, and hence finitely many
possible Frobenius--Perron dimensions.

By Lemma~\ref{lec21-image-dimension}, the fusion image of
$\mathrm{Fr}:\mathcal C\to\mathcal C^{(1)}\boxtimes\mathsf{Ver}_p$ has
dimension at most that of $\mathcal C$. If it had smaller dimension,
minimality would give a symmetric tensor functor from that image to
$\mathsf{Ver}_p$, and composition would contradict the choice of
$\mathcal C$. Consequently equality holds and $\mathrm{Fr}$ is fully
faithful.

If its type is $\Vect$, or $\SVec$ when $p>2$, the lemmas of
Lecture~\ref{lec21} make $\mathcal C$ nondegenerate. Proposition
\ref{lec21-nondegenerate-fiber} then supplies a functor to $\Vect$ or
$\SVec$, followed by the standard inclusion in $\mathsf{Ver}_p$.
This is a contradiction. These cases include every possibility for
$p=2,3$.

Thus $p>3$, and the type must be $\mathsf{Ver}_p^+$ or $\mathsf{Ver}_p$.
By \eqref{lec22-cs-nondegenerate}, the category $\mathcal C_s$ has a super
fiber functor. Compose
\[
 \mathcal C\xrightarrow{\mathrm{Fr}}\mathrm{Fr}(\mathcal C)
 \simeq\mathcal C_s\boxtimes\mathsf{Ver}_p^+
 \longrightarrow\SVec\boxtimes\mathsf{Ver}_p^+
 \simeq\mathsf{Ver}_p.
\]
This is the required functor, contradicting the counterexample. Exactness
is automatic for an additive functor with semisimple source.
\end{proof}

\begin{remark}
This proves the existence statement in \cite[Theorem 1.5]{OstrikVerlinde}.
It does not say that $\mathcal C$ itself embeds in $\mathsf{Ver}_p$:
the functor need not be full. Nor does the theorem assert that every
positive-characteristic symmetric tensor category is semisimple.
The generalization beyond fusion categories requires Frobenius exactness,
which is the subject of the next lecture.
\end{remark}
